\uuid{wL9v}
\exo7id{3926}
\titre{exo7 3926}
\auteur{quercia}
\organisation{exo7}
\datecreate{2010-03-11}
\isIndication{false}
\isCorrection{false}
\chapitre{Fonctions circulaires et hyperboliques inverses}
\sousChapitre{Fonctions circulaires inverses}
\cpge{mpsi-mp2i.calcul-algebrique-trigonometrie.d, pcsi.calcul-algebrique-trigonometrie.d, ptsi.calcul-algebrique-trigonometrie.d, mpsi-mp2i.fonctions-valeurs-reelles-complexes.c, pcsi.techniques-calcul-fonctions.c, ptsi.techniques-calcul-fonctions.c, bcpst1.trigonometrie}
\module{Analyse}
\niveau{L1}
\difficulte{}

\contenu{
\begin{enumerate}
    \item \question{Montrer que : $4\arctan\frac 15 - \arctan \frac 1{239} = \frac \pi4$.}
    \item \question{Montrer que : $\arcsin\frac 45 + \arcsin \frac 5{13} + \arcsin \frac {16}{65} = \frac \pi2$.}
\end{enumerate}
}
